\documentclass[11pt,a4paper]{article}

\usepackage[utf8]{inputenc}
\usepackage[T1]{fontenc}
\usepackage[ngerman]{babel}
\usepackage{amsmath,amssymb,amsthm,mathtools}
\usepackage[a4paper,margin=2.5cm]{geometry}
\usepackage[hidelinks]{hyperref}
\usepackage{booktabs}
\usepackage{array}

\title{Feinauflösung natürlicher Zahlen\\[0.35em]
	\large in Semiprime-Normalform\\[0.2em]
	\normalsize EABC-Restklassen, Provenienzauflösung und $A_2$-Diagonalchirurgie}
\author{}
\date{}

\newtheorem{theorem}{Satz}
\newtheorem{lemma}[theorem]{Lemma}
\newtheorem{proposition}[theorem]{Proposition}
\newtheorem{corollary}[theorem]{Korollar}
\theoremstyle{definition}
\newtheorem{definition}[theorem]{Definition}
\theoremstyle{remark}
\newtheorem{remark}[theorem]{Bemerkung}
\newtheorem{example}[theorem]{Beispiel}

\newcommand{\E}{\mathcal{E}}
\newcommand{\R}{\mathbb{R}}
\newcommand{\Z}{\mathbb{Z}}
\newcommand{\N}{\mathbb{N}}
\newcommand{\one}{\mathbf{1}}
\newcommand{\qq}{\mathbf{q}}
\newcommand{\nn}{\mathbf{n}}
\newcommand{\Pan}{P_{\mathrm{aniso}}}
\newcommand{\Evol}{E_{\mathrm{vol}}}
\newcommand{\EDelta}{E_{\Delta}}
\newcommand{\Esurg}{E_{\mathrm{surg}}}
\newcommand{\forget}{\operatorname{forget}}

\begin{document}
	\maketitle
	
	\begin{abstract}
		\noindent
		Wir entwickeln eine \emph{Feinauflösung} für natürliche Zahlen, die als Produkt
		zweier Primfaktoren $p,q>3$ (Semiprime) vorliegen und in den EABC-Restklassen
		modulo~$12$ liegen. Die Kleinsche Vierergruppe
		$\E=\{E,A,B,C\}\cong C_2\times C_2$ liefert die grobe Multiplikationsalgebra.
		Darüber wird eine provenance-erhaltende Auflösung der Diagonalfaser
		$\mu^{-1}(E)=\Delta_V$ konstruiert, welche die drei Selbstkopplungskanäle
		$E_A,E_B,E_C$ vom Volumenabschluss $\Evol=ABC$ trennt.
		Auf dem Kanalvektor $\qq=(n_A^2,n_B^2,n_C^2)^\top$ wirkt der anisotrope
		Projektor $\Pan=I_3-\tfrac13\one\one^\top$; die chirurgische Energie
		$\Esurg(\qq)=\qq^\top\Pan\qq$ misst die Klassenanisotropie im reellen
		Wurzelraum $A_2^\R$. Die Identifikation mit der Tensorrestauration
		$M_{\mathrm{eff}}=24I_3+D_E^\circ$ bleibt als offene Brücke markiert.
	\end{abstract}
	
	\tableofcontents
	\vspace{1em}
	\hrule
	\vspace{1em}
	
	%==============================================================================
	\section{Einleitung und epistemischer Rahmen}
	\label{sec:intro}
	%==============================================================================
	
	Eine natürliche Zahl $N$ heißt \emph{Semiprime}, wenn
	$N=pq$ mit Primzahlen $p,q$ (nicht notwendig verschieden).
	Für $p,q>3$ liegen die Restklassen in
	\[
	[p]_{12},[q]_{12}\in\{1,5,7,11\},
	\]
	welche wir mit den Symbolen $E,A,B,C$ der EABC-Algebra bezeichnen.
	Die grobe Restklassenmultiplikation kollabiert alle diagonalen Produkte
	$x\cdot x$ auf dasselbe Einselement~$E$. Die \emph{Feinauflösung} besteht darin,
	diese Kollision wieder nach Provenienz zu differenzieren und die resultierende
	anisotropie auf dem Kanalvektor $\qq$ linear zu messen.
	
	Die Architektur ist dreistufig:
	
	\begin{center}
		\begin{tabular}{@{}llp{6.2cm}@{}}
			\toprule
			\textbf{Ebene} & \textbf{Gegenstand} & \textbf{Aussage} \\
			\midrule
			1 & Restklassenalgebra &
			$A^2=B^2=C^2=ABC=E$ in $\E\cong C_2\times C_2$ \\
			2 & Provenienzauflösung &
			$E_A,E_B,E_C,\Evol$ typologisch paarweise verschieden \\
			3 & Lineare Chirurgie &
			$\mathcal S_E(\Evol)=\Evol$,\;
			$\mathcal S_E\big|_{\mathrm{span}\{E_A,E_B,E_C\}}=\Pan$ \\
			\bottomrule
		\end{tabular}
	\end{center}
	
	\begin{remark}[Epistemischer Status]
		\label{rem:status}
		Die Ebenen~1 und~2 sind endlich formal verifizierbar (Lean~4).
		Ebene~3 ist symbolisch geschlossen.
		Die Brücke
		$\Esurg(\qq)\Longrightarrow D_E^\circ\Longrightarrow M_{\mathrm{eff}}=24I_3+D_E^\circ$
		ist ein \emph{Open Non-Claim}.
	\end{remark}
	
	%==============================================================================
	\section{Ebene 1: EABC-Restklassenalgebra}
	\label{sec:eabc}
	%==============================================================================
	
	\begin{definition}[EABC-Klassen]
		\label{def:eabc}
		Die vier invertierbaren Restklassen modulo~$12$ werden bezeichnet durch
		\[
		E\equiv 1,\qquad A\equiv 5,\qquad B\equiv 7,\qquad C\equiv 11\pmod{12}.
		\]
		Wir setzen $\E:=\{E,A,B,C\}$.
	\end{definition}
	
	\begin{definition}[EABC-Multiplikation]
		\label{def:mul}
		Die Multiplikation $\mu:\E\times\E\to\E$ ist die Restklassenmultiplikation
		modulo~$12$ und lautet explizit:
		\[
		\begin{array}{c|cccc}
			\cdot & E & A & B & C \\ \hline
			E & E & A & B & C \\
			A & A & E & C & B \\
			B & B & C & E & A \\
			C & C & B & A & E
		\end{array}
		\]
	\end{definition}
	
	\begin{proposition}[Kleinsche Vierergruppe]
		\label{prop:klein}
		$(\E,\mu)$ ist isomorph zu $C_2\times C_2$. Insbesondere gilt
		$x^2=E$ für alle $x\in\E$, und $\mu$ ist kommutativ und assoziativ.
	\end{proposition}
	
	\begin{proof}
		Unmittelbar aus der Multiplikationstafel; jedes Element ist selbstinvers.
	\end{proof}
	
	\begin{definition}[Nichttrivialer Träger]
		\label{def:V}
		Der nichttriviale Träger ist
		$V:=\{A,B,C\}=\E\setminus\{E\}$.
	\end{definition}
	
	\begin{lemma}[Diagonales und gemischtes Produkt]
		\label{lem:diag-mixed}
		Für $x,y\in V$ gilt:
		\begin{enumerate}
			\item \emph{Diagonaltyp:} falls $x=y$, so $\mu(x,y)=E$;
			\item \emph{Gemischter Typ:} falls $x\neq y$, so existiert ein eindeutig
			bestimmtes $z\in V\setminus\{x,y\}$ mit $\mu(x,y)=z$.
		\end{enumerate}
		Explizit: $AB=C$, $BC=A$, $CA=B$.
	\end{lemma}
	
	\begin{proof}
		Fallunterscheidung an der Tafel; bei $x\neq y$ ist
		$\{x,y,z\}=\{A,B,C\}$.
	\end{proof}
	
	\begin{definition}[Semiprime-Normalform]
		\label{def:semiprime-nf}
		Ein Semiprime $N=pq$ mit $p,q>3$ liegt in Semiprime-Normalform vor, wenn
		$[p]_{12},[q]_{12}\in V$ und
		\[
		\mu\bigl([p]_{12},[q]_{12}\bigr)
		=
		\begin{cases}
			E, & [p]_{12}=[q]_{12},\\
			z, & \{[p]_{12},[q]_{12},z\}=\{A,B,C\}.
		\end{cases}
		\]
	\end{definition}
	
	\begin{lemma}[Tripelabschluss]
		\label{lem:ABC}
		Für paarweise verschiedene $x,y,z\in V$ gilt
		$\mu(\mu(x,y),z)=E$. Insbesondere $ABC=E$ in $\E$.
	\end{lemma}
	
	\begin{proof}
		Wegen Lemma~\ref{lem:diag-mixed} ist $\mu(x,y)=z'$, wobei $\{x,y,z'\}=\{A,B,C\}$.
		Bei $\{x,y,z\}=\{A,B,C\}$ folgt $z'=z$, also $\mu(z,z)=E$.
	\end{proof}
	
	\begin{remark}
		\label{rem:same-E}
		In der Restklassenalgebra sind $A^2=E$ und $ABC=E$ \emph{dasselbe}
		Gruppenelement. Die Unterscheidung $\EDelta\neq\Evol$ erfordert Ebene~2.
	\end{remark}
	
	%==============================================================================
	\section{Diagonalfaser und Semiprime-Struktur}
	\label{sec:fiber}
	%==============================================================================
	
	\begin{definition}[Multiplikationsfaser]
		\label{def:fiber}
		Für $w\in\E$ sei $\mu^{-1}(w):=\{(x,y)\in V\times V:\mu(x,y)=w\}$.
	\end{definition}
	
	\begin{theorem}[Diagonalfaser]
		\label{thm:diagonal-fiber}
		Es gilt
		\[
		\mu^{-1}(E)=\Delta_V:=\{(A,A),(B,B),(C,C)\}.
		\]
		Die Fasern der nichttrivialen Ausgänge sind dagegen gemischt:
		\begin{align*}
			\mu^{-1}(A)&=\{(B,C),(C,B)\},\\
			\mu^{-1}(B)&=\{(A,C),(C,A)\},\\
			\mu^{-1}(C)&=\{(A,B),(B,A)\}.
		\end{align*}
	\end{theorem}
	
	\begin{proof}
		Direkt aus Lemma~\ref{lem:diag-mixed}.
	\end{proof}
	
	\begin{corollary}[Strukturelle Asymmetrie]
		\label{cor:asymmetry}
		$A,B,C$ entstehen aus gemischten Paaren; $E$ entsteht auf $V\times V$
		ausschließlich aus Selbstpaarungen.
	\end{corollary}
	
	%==============================================================================
	\section{Ebene 2: Provenienzauflösung}
	\label{sec:provenance}
	%==============================================================================
	
	\begin{definition}[Provenienzmarker]
		\label{def:prov}
		Die Menge der Provenienzmarker über dem Restklassenbild~$E$ ist
		\[
		\mathrm{EProvenance}
		:=
		\{E_A,E_B,E_C,\Evol\}.
		\]
	\end{definition}
	
	\begin{definition}[Vergessensabbildung]
		\label{def:forget}
		Die Vergessensabbildung
		$\forget:\mathrm{EProvenance}\to\E$
		ist durch
		$\forget(E_A)=\forget(E_B)=\forget(E_C)=\forget(\Evol)=E$
		definiert.
	\end{definition}
	
	\begin{definition}[Aufgelöster Ausgang]
		\label{def:resolved}
		Der aufgelöste Ausgangsraum ist die disjunkte Summe
		\[
		\mathrm{ResolvedOutput}
		:=
		V\;\sqcup\;\mathrm{EProvenance},
		\]
		typologisch geschrieben als
		$\texttt{nontrivial}\oplus\texttt{prov}$.
	\end{definition}
	
	\begin{definition}[Aufgelöstes Produkt]
		\label{def:resolved-mul}
		Das provenance-erhaltende Produkt
		$\widetilde\mu:V\times V\to\mathrm{ResolvedOutput}$
		ist gegeben durch
		\[
		\widetilde\mu(x,y)
		=
		\begin{cases}
			E_x, & x=y,\\
			z, & x\neq y,\;\{x,y,z\}=\{A,B,C\}.
		\end{cases}
		\]
		Explizit:
		\[
		\begin{array}{c|ccc}
			\widetilde\mu & A & B & C \\ \hline
			A & E_A & C & B \\
			B & C & E_B & A \\
			C & B & A & E_C
		\end{array}
		\]
	\end{definition}
	
	\begin{definition}[Projektion]
		\label{def:project}
		Die Projektion $\pi:\mathrm{ResolvedOutput}\to\E$ ist
		$\pi|_V=\mathrm{id}_V$ und
		$\pi|_{\mathrm{EProvenance}}=\forget$.
	\end{definition}
	
	\begin{theorem}[Kommutatives Quadrat]
		\label{thm:commutes}
		Für alle $x,y\in V$ gilt
		\[
		\pi\circ\widetilde\mu=\mu\Big|_{V\times V}.
		\]
	\end{theorem}
	
	\begin{proof}
		Auf der Diagonale: $\pi(E_x)=E=\mu(x,x)$.
		Off-Diagonal: $\pi(z)=z=\mu(x,y)$.
	\end{proof}
	
	\begin{theorem}[Provenienztrennung]
		\label{thm:provenance}
		Es gilt
		\[
		\Evol\notin\{E_A,E_B,E_C\}
		\quad\text{und}\quad
		\forget(\Evol)=\forget(E_x)=E\quad(x\in V).
		\]
		Somit sind $\EDelta$ (diagonale Selbstkopplung) und $\Evol$ (Tripelabschluss)
		verschiedene Provenienzklassen mit gleichem Restklassenbild.
	\end{theorem}
	
	\begin{definition}[Annotiertes Tripel]
		\label{def:triple}
		Die annotierte Auswertung $\tau:V^3\to\mathrm{ResolvedOutput}$ ist
		\[
		\tau(x,y,z)
		=
		\begin{cases}
			\Evol, & x,y,z\text{ paarweise verschieden},\\
			x, & x=y=z,\\
			\text{Trägerklasse}, & \text{genau zwei gleich}
		\end{cases}
		\]
		(bei genau zwei gleichen Faktoren liefert die Halsheilung $x^2\cdot y=y$
		wieder eine Klasse in~$V$).
	\end{definition}
	
	\begin{lemma}[Halsheilung]
		\label{lem:neck}
		Für $x,y\in V$ gilt in $\E$:
		$\mu(\mu(x,x),y)=y$.
		Im annotierten Sinn wird der diagonale Zwischenmarker durch den transversalen
		Faktor wieder in den Träger zurückgeführt.
	\end{lemma}
	
	\begin{proof}
		$\mu(x,x)=E$ und $\mu(E,y)=y$.
	\end{proof}
	
	\begin{proposition}[Volumenabschluss]
		\label{prop:volume}
		Es gilt $\tau(A,B,C)=\Evol$ und
		$\widetilde\mu(x,x)\neq\Evol$ für alle $x\in V$.
	\end{proposition}
	
	%==============================================================================
	\section{Kanalvektoren und Zählkonventionen}
	\label{sec:channels}
	%==============================================================================
	
	\begin{definition}[Faktorhäufigkeitsvektor]
		\label{def:n}
		Seien $n_A,n_B,n_C\in\R_{\ge 0}$ die Anzahlen der Primfaktoren in den Klassen
		$A,B,C$. Der Häufigkeitsvektor ist
		$\nn=(n_A,n_B,n_C)^\top$.
	\end{definition}
	
	\begin{definition}[Paarzählkonvention]
		\label{def:convention}
		Die diagonale Semiprimekanalmasse $q_x$ hängt von der Multimengenkonvention ab:
		\begin{align*}
			q_x&=n_x^2
			&&\text{(geordnete Paare mit Wiederholung)},\\
			q_x&=n_x(n_x-1)
			&&\text{(geordnete Paare, $p\neq q$)},\\
			q_x&=\binom{n_x}{2}
			&&\text{(ungeordnete Paare, $p\neq q$)},\\
			q_x&=\binom{n_x+1}{2}
			&&\text{(ungeordnet inkl.\ Primquadrate)}.
		\end{align*}
		Im Folgenden ist die Standardkonvention $q_x=n_x^2$.
	\end{definition}
	
	\begin{definition}[Kanalvektor]
		\label{def:q}
		Der diagonale Semiprimekanalvektor ist
		$\qq=(q_A,q_B,q_C)^\top$ mit $q_x=n_x^2$.
	\end{definition}
	
	\begin{remark}
		\label{rem:n-vs-q}
		Die Größe $\nn^\top\Pan\nn$ misst Häufigkeitsanisotropie;
		$\qq^\top\Pan\qq$ misst die Anisotropie der tatsächlich erzeugten
		Selbstkopplungskanäle. Die Chirurgie operiert auf~$\qq$.
	\end{remark}
	
	%==============================================================================
	\section{Ebene 3: Anisotroper Projektor und $A_2$-Energie}
	\label{sec:paniso}
	%==============================================================================
	
	\begin{definition}[Anisotroper Projektor]
		\label{def:paniso}
		Auf $\R^3$ sei
		\[
		\Pan
		:=
		I_3-\frac13\one\one^\top,
		\qquad
		\one=(1,1,1)^\top.
		\]
	\end{definition}
	
	\begin{definition}[Reeller und ganzzahliger $A_2$-Raum]
		\label{def:a2}
		\[
		A_2^\R
		:=
		\{(a,b,c)\in\R^3:a+b+c=0\},
		\qquad
		A_2
		:=
		\{(a,b,c)\in\Z^3:a+b+c=0\}.
		\]
	\end{definition}
	
	\begin{lemma}[Projektoraxiome]
		\label{lem:projector}
		Es gilt:
		\begin{enumerate}
			\item Idempotenz: $\Pan^2=\Pan$;
			\item Kern: $\Pan\one=0$;
			\item Bild: $\operatorname{range}(\Pan)=A_2^\R$.
		\end{enumerate}
	\end{lemma}
	
	\begin{proof}
		(1)~folgt aus $\one^\top\one=3$.
		(2)~ist $\Pan\one=\one-\one=0$.
		(3)~: $\Pan v$ ist spurfrei; ist $v$ spurfrei, so $\Pan v=v$.
	\end{proof}
	
	\begin{remark}
		\label{rem:lattice}
		Für ganzzahliges $\nn$ liegt $\Pan\nn$ im reellen Raum $A_2^\R$, wegen des
		Faktors $\tfrac13$ jedoch nicht notwendig im Gitter~$A_2$.
		Ganzzahlige Defekte liefert etwa
		$(n_A-n_B,\,n_B-n_C,\,n_C-n_A)^\top\in A_2$.
	\end{remark}
	
	\begin{definition}[Chirurgische Energie]
		\label{def:esurg}
		\[
		\Esurg(\qq)
		:=
		\qq^\top\Pan\qq.
		\]
	\end{definition}
	
	\begin{theorem}[$A_2$-Energieidentität]
		\label{thm:energy}
		Für jeden Kanalvektor $\qq\in\R^3$ gilt
		\[
		\Esurg(\qq)
		=
		\qq^\top\Pan\qq
		=
		\frac13\Bigl[
		(q_A-q_B)^2+(q_B-q_C)^2+(q_C-q_A)^2
		\Bigr].
		\]
		Insbesondere $\Esurg(\qq)\ge 0$, und
		$\Esurg(\qq)=0$ genau dann, wenn $q_A=q_B=q_C$.
	\end{theorem}
	
	\begin{proof}
		Mit $\bar q=\tfrac13(q_A+q_B+q_C)$ und $\delta_x=q_x-\bar q$ ist
		$\Pan\qq=\boldsymbol\delta$ und
		$\qq^\top\Pan\qq=\|\boldsymbol\delta\|^2$.
		Die Identität
		$\sum_x\delta_x^2=\tfrac13\sum_{\mathrm{zyk}}(q_x-q_y)^2$
		folgt durch direkte Expansion unter $\delta_A+\delta_B+\delta_C=0$.
	\end{proof}
	
	\begin{corollary}[Zerlegung der Diagonalmasse]
		\label{cor:decomp}
		Mit $M_E=\|\qq\|_2^2$ (geordnete Paare) gilt
		\[
		M_E
		=
		3\bar q^2+\Esurg(\qq).
		\]
		Der Term $3\bar q^2$ ist der \emph{isotrope Diagonalanteil};
		$\Esurg$ ist der anisotrope Rest.
	\end{corollary}
	
	%==============================================================================
	\section{Chirurgischer Operator mit Volumenschutz}
	\label{sec:surgery}
	%==============================================================================
	
	\begin{definition}[Erweiterter Diagonalraum]
		\label{def:extended}
		Der erweiterte Raum der $E$-Provenienzen ist
		\[
		\widetilde V_E
		:=
		\operatorname{span}\{E_A,E_B,E_C\}
		\;\oplus\;
		\operatorname{span}\{\Evol\}
		\;\cong\;
		\R^3\oplus\R.
		\]
	\end{definition}
	
	\begin{definition}[Chirurgischer Operator]
		\label{def:SE}
		Der Operator $\mathcal S_E:\widetilde V_E\to\widetilde V_E$ ist definiert durch
		\[
		\mathcal S_E\big|_{\operatorname{span}\{E_A,E_B,E_C\}}
		=
		\Pan,
		\qquad
		\mathcal S_E(\Evol)=\Evol.
		\]
		Äquivalent in Koordinaten $(e_A,e_B,e_C;e_{\mathrm{vol}})$:
		\[
		\mathcal S_E(e_A,e_B,e_C;e_{\mathrm{vol}})
		=
		\bigl(e_A-\bar e,\,e_B-\bar e,\,e_C-\bar e;\,e_{\mathrm{vol}}\bigr),
		\]
		wobei $\bar e=\tfrac13(e_A+e_B+e_C)$.
	\end{definition}
	
	\begin{theorem}[Semiprime-Diagonalchirurgie]
		\label{thm:main}
		Es gilt:
		\begin{enumerate}
			\item Diagonalfaser: $\mu^{-1}(E)=\Delta_V$ auf $V\times V$.
			\item Die Auflösung $\widetilde\mu$ trennt $E_A,E_B,E_C$.
			\item $\Evol$ liegt nicht in der Diagonalfaser und ist unter $\mathcal S_E$
			invariant.
			\item $\mathcal S_E$ vernichtet genau den isotropen Diagonalanteil:
			$\mathcal S_E(E_A+E_B+E_C)=0$ im Diagonalsektor.
			\item Die Restenergie $\Esurg(\qq)$ verschwindet genau bei
			Kanalisotropie $q_A=q_B=q_C$.
		\end{enumerate}
	\end{theorem}
	
	\begin{proof}
		(1)--(2) aus den Theoremen~\ref{thm:diagonal-fiber} und~\ref{thm:commutes}.
		(3)--(4) aus Definition~\ref{def:SE} und Lemma~\ref{lem:projector}.
		(5) aus Theorem~\ref{thm:energy}.
	\end{proof}
	
	\begin{corollary}[Leitprinzip]
		\label{cor:principle}
		Die Chirurgie wirkt auf den aufgelösten diagonalen Kanalvektor~$\qq$,
		nicht auf das Gruppenelement $E\in\E$.
	\end{corollary}
	
	%==============================================================================
	\section{Offene Brücke zum $24I_3$-Tensorsektor}
	\label{sec:bridge}
	%==============================================================================
	
	\begin{definition}[Diagonale Einbettung]
		\label{def:embed}
		Die lineare Einbettung $\iota:\R^3\to\mathrm{Sym}_3(\R)$ ist
		$\iota(a,b,c)=\operatorname{diag}(a,b,c)$.
	\end{definition}
	
	\begin{definition}[Defekttensor-Kandidat]
		\label{def:DE}
		Für einen Kanalvektor $\qq$ setzen wir
		\[
		D_E^\circ(\qq)
		:=
		\iota(\Pan\qq)
		=
		\operatorname{diag}(\delta_A,\delta_B,\delta_C),
		\]
		wobei $\boldsymbol\delta=\Pan\qq$, und
		\[
		M_{\mathrm{eff}}^{\mathrm{cand}}(\qq)
		:=
		24I_3+D_E^\circ(\qq).
		\]
	\end{definition}
	
	\begin{proposition}[Schnittstelleneigenschaften]
		\label{prop:interface}
		Es gilt $\operatorname{tr}D_E^\circ(\qq)=0$ und
		\[
		\bigl\|D_E^\circ(\qq)\bigr\|_F^2
		=
		\Esurg(\qq).
		\]
	\end{proposition}
	
	\begin{proof}
		Spurfreiheit folgt aus $\Pan\qq\in A_2^\R$.
		Die Frobeniusnorm diagonaler Matrizen reduziert sich auf
		$\|\Pan\qq\|_2^2=\qq^\top\Pan\qq$.
	\end{proof}
	
	\begin{remark}[Open Non-Claim]
		\label{rem:open}
		Die Identifikation von $M_{\mathrm{eff}}^{\mathrm{cand}}$ mit der
		Restaurationssemantik der EABC-Operatoralgebra
		($M_{\mathrm{eff}}(R^*(K))=24I_3$) ist \emph{nicht} bewiesen.
		Strukturelle Alternativen zur Einbettung sind:
		\begin{enumerate}
			\item linear: $\iota\circ\Pan$ (erhält $\dim A_2^\R=2$);
			\item nichtlinear: spurfreie $\operatorname{diag}(\delta_A^2,\delta_B^2,\delta_C^2)$;
			\item Rang-1-Kollaps $w\,vv^\top$ (nur eindimensional; Dimension geht verloren).
		\end{enumerate}
	\end{remark}
	
	%==============================================================================
	\section{Zusammenfassung}
	\label{sec:summary}
	%==============================================================================
	
	Die Feinauflösung Semiprimer in EABC-Normalform trennt drei Ebenen:
	
	\begin{enumerate}
		\item In $\E$ kollabieren Diagonale und Volumen auf dasselbe $E$.
		\item Die Provenienzauflösung stellt $E_A,E_B,E_C$ und $\Evol$ als
		verschiedene Marker wieder her.
		\item Die lineare Chirurgie $\Pan$ kontrahiert den isotropen Diagonalanteil
		auf dem Kanalvektor $\qq$ und schützt $\Evol$; die Energie
		$\Esurg(\qq)$ ist die exakte $A_2$-Varianz der Selbstkopplungskanäle.
	\end{enumerate}
	
	Die einzige offene Schicht ist die kanonische Einbettung dieses Defekts in den
	bestehenden $3\times 3$-Tensorsektor der $24I_3$-Restauration.
	
	\begin{thebibliography}{9}
		
		\bibitem{eabc-renorm}
		T.~Hoffbauer,
		\emph{EABC-Operatoralgebra und Renormierungskette},
		Manuskript und Lean-4-Formalisierung,
		\texttt{eabc-renorm}.
		
		\bibitem{kepler-hurwitz}
		KeplerHurwitz/EABC,
		\emph{SemiprimeNormalForm}, \emph{ProvenanceComposition},
		\emph{SemiprimeTensorSurgery}, \emph{Semiprime24Bridge},
		Lean~4.29-Module.
		
	\end{thebibliography}
	
\end{document}
